% GENERATED FILE. Edit canonical lessons, then run npm run generate:books.
% canonical-source-hash: fnv1a64:813fe983f8aac3d0
% canonical-lessons: TA-C393-arikkai, TA-C393-kanakketuppu, TA-C393-kantari, TA-C393-cumar, TA-C393-perumpalum

\chapter{Facts and Figures}
\label{ch:ta-facts-and-figures}
\hlchaptermodality{393}

\begin{chapteropening}
\textbf{By the end of this chapter:} \emph{I can talk about reports, surveys and rough numbers.}

Complete the last lesson of chapter 393: I can talk about reports, surveys and rough numbers.
\end{chapteropening}

\section[\texorpdfstring{\ta{அறிக்கை} (aṟikkai)}{அறிக்கை (aṟikkai)}]{\ta{அறிக்கை} (aṟikkai) — a report; a statement}

\label{lesson:TA-C393-arikkai}



\begin{quote}
Before the new one: say the Tamil for indirectly, then the Tamil for to monitor.
\end{quote}

\subsection*{You'll want to know: \ta{அறிக்கை}}
\textbf{\ta{அறிக்கை}} — \emph{aṟikkai} — \textquotedblleft{}a report; a statement\textquotedblright{}.

\begin{quote}
\ta{அறிக்கை} \ta{வந்தால்} \ta{உண்மை} \ta{தெரியும்}.
\end{quote}

\begin{quote}
\emph{aṟikkai vantāl uṇmai teriyum.}
\end{quote}

\begin{quote}
If the report comes out, the truth will be known.
\end{quote}

\ta{அறிக்கை} \ta{வெளியிடு}, to issue a statement.

\begin{cousinweb}
\textbf{\ta{அறிக்கை}} (\emph{aṟikkai}) is native Dravidian: DEDR 314 lists Tamil \emph{aṟi}, \textquotedblleft{}to know\textquotedblright{}, and \emph{aṟikkai}, \textquotedblleft{}notice\textquotedblright{}: what makes a thing known.
\end{cousinweb}

\begin{grammarlens}[title={What you've built}]
One more word for saying what happens if.
\end{grammarlens}

\subsection*{Your turn}
\begin{itemize}
\raggedright
  \item \emph{Say it:} \emph{aṟikkai}
  \item \emph{Say it:} \emph{aṟikkai}, once more
  \item \emph{Say it:} \emph{aṟikkai vantāl uṇmai teriyum.} — If the report comes out, the truth will be known.
  \item \emph{Recall:} say the Tamil for indirectly, then the Tamil for to monitor, then say \emph{aṟikkai} again
\end{itemize}

\begin{tcolorbox}[breakable,colback=teal!4,colframe=teal!35!black,title={Before you move on}]
What does \ta{அறிக்கை} mean? (\textquotedblleft{}A report; a statement\textquotedblright{}.) Say it once more.
\end{tcolorbox}
\section[\texorpdfstring{\ta{கணக்கெடுப்பு} (kaṇakkeṭuppu)}{கணக்கெடுப்பு (kaṇakkeṭuppu)}]{\ta{கணக்கெடுப்பு} (kaṇakkeṭuppu) — a census; a survey}

\label{lesson:TA-C393-kanakketuppu}



\begin{quote}
Before the new one: say the Tamil for to monitor, then the Tamil for a report.
\end{quote}

\subsection*{You'll want to know: \ta{கணக்கெடுப்பு}}
\textbf{\ta{கணக்கெடுப்பு}} — \emph{kaṇakkeṭuppu} — \textquotedblleft{}a census; a survey\textquotedblright{}.

\begin{quote}
\ta{பத்து} \ta{ஆண்டுகளுக்கு} \ta{ஒருமுறை} \ta{மக்கள்தொகைக்} \ta{கணக்கெடுப்பு} \ta{நடக்கிறது}.
\end{quote}

\begin{quote}
\emph{pattu āṇṭukaḷukku orumuṟai makkaḷtokaik kaṇakkeṭuppu naṭakkiṟatu.}
\end{quote}

\begin{quote}
A population census is held once every ten years.
\end{quote}

Literally a taking of the count: \ta{எடு}, to take, is inside it.

\begin{cousinweb}
\textbf{\ta{கணக்கெடுப்பு}} (\emph{kaṇakkeṭuppu}) is built from \emph{kaṇakku}, \textquotedblleft{}an account, a count\textquotedblright{}, and native \emph{eṭuppu}, \textquotedblleft{}a taking up\textquotedblright{}, from \emph{eṭu}, \textquotedblleft{}to take\textquotedblright{} (DEDR 851).
\end{cousinweb}

\begin{grammarlens}[title={What you've built}]
One more word for saying what happens if.
\end{grammarlens}

\subsection*{Your turn}
\begin{itemize}
\raggedright
  \item \emph{Say it:} \emph{kaṇakkeṭuppu}
  \item \emph{Say it:} \emph{kaṇakkeṭuppu}, once more
  \item \emph{Say it:} \emph{pattu āṇṭukaḷukku orumuṟai makkaḷtokaik kaṇakkeṭuppu naṭakkiṟatu.} — A population census is held once every ten years.
  \item \emph{Recall:} say the Tamil for to monitor, then the Tamil for a report, then say \emph{kaṇakkeṭuppu} again
\end{itemize}

\begin{tcolorbox}[breakable,colback=teal!4,colframe=teal!35!black,title={Before you move on}]
What does \ta{கணக்கெடுப்பு} mean? (\textquotedblleft{}A census; a survey\textquotedblright{}.) Say it once more.
\end{tcolorbox}
\section[\texorpdfstring{\ta{கண்டறி} (kaṇṭaṟi)}{கண்டறி (kaṇṭaṟi)}]{\ta{கண்டறி} (kaṇṭaṟi) — to find out; to discover}

\label{lesson:TA-C393-kantari}



\begin{quote}
Before the new one: say the Tamil for a report, then the Tamil for a census.
\end{quote}

\subsection*{You'll want to know: \ta{கண்டறி}}
\textbf{\ta{கண்டறி}} — \emph{kaṇṭaṟi} — \textquotedblleft{}to find out; to discover\textquotedblright{}.

\begin{quote}
\ta{காரணத்தைக்} \ta{கண்டறிந்தால்} \ta{தீர்வு} \ta{கிடைக்கும்}.
\end{quote}

\begin{quote}
\emph{kāraṇattaik kaṇṭaṟintāl tīrvu kiṭaikkum.}
\end{quote}

\begin{quote}
If we find out the cause, we will find a solution.
\end{quote}

Past (I …): \ta{கண்டறிந்தேன்}, \emph{kaṇṭaṟintēṉ}.

What is found takes -\ta{ஐ}. A more formal word than \ta{கண்டுபிடி}.

\begin{cousinweb}
\textbf{\ta{கண்டறி}} (\emph{kaṇṭaṟi}) is native: \emph{kaṇṭu}, \textquotedblleft{}having seen\textquotedblright{}, from \emph{kāṇ} (DEDR 1443), and \emph{aṟi}, \textquotedblleft{}to know\textquotedblright{} (DEDR 314): to see and come to know.
\end{cousinweb}

\begin{grammarlens}[title={What you've built}]
One more word for saying what happens if.
\end{grammarlens}

\subsection*{Your turn}
\begin{itemize}
\raggedright
  \item \emph{Say it:} \emph{kaṇṭaṟi}
  \item \emph{Say it:} \emph{kaṇṭaṟi}, once more
  \item \emph{Say it:} \emph{kāraṇattaik kaṇṭaṟintāl tīrvu kiṭaikkum.} — If we find out the cause, we will find a solution.
  \item \emph{Recall:} say the Tamil for a report, then the Tamil for a census, then say \emph{kaṇṭaṟi} again
\end{itemize}

\begin{tcolorbox}[breakable,colback=teal!4,colframe=teal!35!black,title={Before you move on}]
What does \ta{கண்டறி} mean? (\textquotedblleft{}To find out; to discover\textquotedblright{}.) Say it once more.
\end{tcolorbox}
\section[\texorpdfstring{\ta{சுமார்} (cumār)}{சுமார் (cumār)}]{\ta{சுமார்} (cumār) — roughly; about}

\label{lesson:TA-C393-cumar}



\begin{quote}
Before the new one: say the Tamil for a census, then the Tamil for to find out.
\end{quote}

\subsection*{You'll want to know: \ta{சுமார்}}
\textbf{\ta{சுமார்}} — \emph{cumār} — \textquotedblleft{}roughly; about\textquotedblright{}.

\begin{quote}
\ta{கூட்டத்துக்குச்} \ta{சுமார்} \ta{நூறு} \ta{பேர்} \ta{வந்தார்கள்}.
\end{quote}

\begin{quote}
\emph{kūṭṭattukkuc cumār nūṟu pēr vantārkaḷ.}
\end{quote}

\begin{quote}
About a hundred people came to the meeting.
\end{quote}

Before a number: \ta{சுமார்} \ta{பத்து} \ta{மணிக்கு}, at about ten o'clock. In speech it also means so-so.

\begin{cousinweb}
\textbf{\ta{சுமார்}} (\emph{cumār}) is usually traced to Persian \emph{shumār}, \textquotedblleft{}counting, number\textquotedblright{}, which reached Tamil through Hindustani: a figure by rough count.
\end{cousinweb}

\begin{grammarlens}[title={What you've built}]
One more word for saying what happens if.
\end{grammarlens}

\subsection*{Your turn}
\begin{itemize}
\raggedright
  \item \emph{Say it:} \emph{cumār}
  \item \emph{Say it:} \emph{cumār}, once more
  \item \emph{Say it:} \emph{kūṭṭattukkuc cumār nūṟu pēr vantārkaḷ.} — About a hundred people came to the meeting.
  \item \emph{Recall:} say the Tamil for a census, then the Tamil for to find out, then say \emph{cumār} again
\end{itemize}

\begin{tcolorbox}[breakable,colback=teal!4,colframe=teal!35!black,title={Before you move on}]
What does \ta{சுமார்} mean? (\textquotedblleft{}Roughly; about\textquotedblright{}.) Say it once more.
\end{tcolorbox}
\section[\texorpdfstring{\ta{பெரும்பாலும்} (perumpālum)}{பெரும்பாலும் (perumpālum)}]{\ta{பெரும்பாலும்} (perumpālum) — mostly; usually}

\label{lesson:TA-C393-perumpalum}



\begin{quote}
Before the new one: say the Tamil for to find out, then the Tamil for roughly.
\end{quote}

\subsection*{You'll want to know: \ta{பெரும்பாலும்}}
\textbf{\ta{பெரும்பாலும்}} — \emph{perumpālum} — \textquotedblleft{}mostly; usually\textquotedblright{}.

\begin{quote}
\ta{மழை} \ta{பெய்தால்} \ta{பெரும்பாலும்} \ta{நான்} \ta{வீட்டிலேயே} \ta{இருப்பேன்}.
\end{quote}

\begin{quote}
\emph{maḻai peytāl perumpālum nāṉ vīṭṭilēyē iruppēṉ.}
\end{quote}

\begin{quote}
If it rains, I mostly stay at home.
\end{quote}

\ta{பெரும்}, great + \ta{பால்}, part + -\ta{உம்}: for the greater part.

\begin{cousinweb}
\textbf{\ta{பெரும்பாலும்}} (\emph{perumpālum}) is native: \emph{perum}, \textquotedblleft{}great\textquotedblright{} (DEDR 4411), \emph{pāl}, \textquotedblleft{}part, portion\textquotedblright{} (DEDR 4097), and the particle -\ta{உம்}: for the greater part.
\end{cousinweb}

\begin{grammarlens}[title={What you've built}]
One more word for saying what happens if.
\end{grammarlens}

\subsection*{Your turn}
\begin{itemize}
\raggedright
  \item \emph{Say it:} \emph{perumpālum}
  \item \emph{Say it:} \emph{perumpālum}, once more
  \item \emph{Say it:} \emph{maḻai peytāl perumpālum nāṉ vīṭṭilēyē iruppēṉ.} — If it rains, I mostly stay at home.
  \item \emph{Recall:} say the Tamil for to find out, then the Tamil for roughly, then say \emph{perumpālum} again
\end{itemize}

\begin{tcolorbox}[breakable,colback=teal!4,colframe=teal!35!black,title={Before you move on}]
What does \ta{பெரும்பாலும்} mean? (\textquotedblleft{}Mostly; usually\textquotedblright{}.) Say it once more.
\end{tcolorbox}
